Esta tesis estudia el problema de construir un \zkBridge para \Cardano. En términos generales, un \Bridge ---noción introducida en la Sección~\ref{subsec:bridges-que-es-un-bridge}--- busca transferir valor o información entre dos cadenas sin exigir a sus usuarios confianza plena en un intermediario centralizado. En nuestro caso, el objetivo es más específico: diseñar una arquitectura que permita tomar un evento ocurrido en una cadena de origen, autenticarlo criptográficamente y habilitar una acción correspondiente en una cadena destino, todo ello respetando las restricciones reales del entorno de ejecución de \Cardano.

El desafío no es solamente criptográfico. Un \zkBridge práctico debe combinar, al mismo tiempo, un mecanismo de autenticación del estado remoto, una representación compacta de la evidencia relevante, una estrategia de verificación \Onchain compatible con límites estrictos de cómputo y tamaño, y un modelo de estado que preserve invariantes como la atomicidad y la prevención de \textit{replay}. En \Cardano, estas exigencias quedan además condicionadas por el modelo \eUTxO, desarrollado en la Sección~\ref{sec:motivation-el-modelo-e-utxo}, y por los límites de ejecución reunidos en la Sección~\ref{sec:motivation-limites-del-protocolo-y-su-importancia-en-un-zkbridge}.

La propuesta desarrollada en este trabajo parte de una decisión central: no intentar verificar directamente dentro de un circuito \ZK todo el consenso nativo de \Cardano, sino apoyarse en \Mithril como mecanismo de certificación compacta del estado relevante, estudiado en el Capítulo~\ref{chap:mithril}. Sus certificados se autentican mediante un esquema de \textit{stake-based threshold multisignatures} (\STM), explicado en la Sección~\ref{sec:mithril-stm-protocol}, cuya verificación \HaloTwo/\Midnight se divide en dos fases según la Sección~\ref{subsec:mithril-a-cardano-particion-de-la-verificacion-en-dos-transacciones}. Para los circuitos auxiliares de pertenencia a \textit{snapshots} y actualización del \SMT anti-replay se utiliza \Groth, presentado en la Sección~\ref{sec:aritmetizacion-y-construccion-de-un-snark-groth}. Sobre esa base, la tesis organiza el \zkBridge en una capa de actualización, que mantiene en la cadena destino un estado autenticado y reutilizable, y una capa de aplicación, que consume esa autenticación para autorizar operaciones concretas del protocolo.

Para seguir el cuerpo principal de la tesis conviene tener familiaridad previa con nociones básicas de blockchains, transacciones, contratos inteligentes y criptografía de clave pública. También se utilizan de manera recurrente herramientas algebraicas y criptográficas como cuerpos finitos, curvas elípticas, funciones de hash, Merkle Trees, argumentos de conocimiento cero, circuitos aritméticos, \ROneCS, \QAP, \KZG, \FiatShamir y \Groth. Esos temas no se desarrollan linealmente dentro del cuerpo principal, para no interrumpir la discusión sobre el diseño del \zkBridge, sino que se recopilan como material de apoyo al final del manuscrito: el Apéndice~\ref{chap:primitivas-criptograficas} cubre las primitivas criptográficas y estructuras autenticadas de base, mientras que el Apéndice~\ref{chap:aritmetizacion-y-construccion-de-un-snark} desarrolla la aritmetización y la construcción de un \SNARK. El Apéndice~\ref{chap:modelo-homogeneo-cadenas-cardano-compatibles}, por su parte, ofrece una formulación más rigurosa del modelo homogéneo adoptado para los roles de origen y destino y de la relación entre \LockTx y \MintTx. Su lectura no es un prerrequisito: cuando el cuerpo principal introduzca esa abstracción, indicará expresamente que el lector puede acudir a dicho apéndice si desea seguir la exposición con mayor formalismo. En cambio, los conceptos específicos de \Cardano, \eUTxO, \Mithril y del flujo implementado se introducen progresivamente en los capítulos centrales.

Con ese marco, el recorrido del manuscrito avanza desde los antecedentes sobre \Bridges y \zkBridge del Capítulo~\ref{chap:bridges-y-zkbridge} hacia las restricciones de \Cardano del Capítulo~\ref{chap:cardano-y-modelo-eutxo}, la certificación con \Mithril del Capítulo~\ref{chap:mithril} y las decisiones arquitectónicas del Capítulo~\ref{chap:decisiones-de-diseno}. El Capítulo~\ref{chap:implementacion-del-zkbridge} concreta luego esas decisiones en una implementación por etapas del camino principal \(\LockTx \rightarrow \MintTx\), cuyos resultados se evalúan en el Capítulo~\ref{chap:resultados}. Los fundamentos necesarios para sostener esa discusión pueden consultarse en los Apéndices~\ref{chap:primitivas-criptograficas} y~\ref{chap:aritmetizacion-y-construccion-de-un-snark}; la formalización opcional del modelo de cadenas se encuentra en el Apéndice~\ref{chap:modelo-homogeneo-cadenas-cardano-compatibles}.
